\section{Setup}\label{sec:setup}

We work on a probability space $(\Omega,\mathcal F,\PP)$. The covariates are a random variable $X\colon\Omega\to\X$, with $(\X,\A)$ a standard Borel space; this guarantees regular conditional distributions and allows mixed vectors of continuous, categorical and missing variables.

\begin{center}
\begin{tabular}{@{}ll@{}}
\toprule
Symbol & Meaning \\
\midrule
$D\in\{0,1\}$ & Channel serving the unit (1 = premium, 0 = basic) \\
$S\in\{0,1\}$ & Membership in the target group, with $\{S=1\}\subseteq\{D=1\}$ \\
$Y(0),\,Y(1)$ & Potential outcomes if served by the basic or the premium channel \\
$Y$ & Observed outcome \\
$P$ & Law of $X$ in the target group: $P(A)=\PP(X\in A\mid S=1)$ \\
$Q$ & Law of $X$ in the donor channel: $Q(A)=\PP(X\in A\mid D=0)$ \\
\bottomrule
\end{tabular}
\end{center}

We assume $\PP(S=1)>0$ and $\PP(D=0)>0$. A typical case is $S=\ind\{D=1,\ s(X)\le c\}$: the target group consists of the premium-channel units whose score $s(X)$ falls below a threshold, for instance the tail with the lowest purchase propensity.

\paragraph{Estimand.} The counterfactual mean of the target group in the basic channel,
\[
\theta=\E\bigl[Y(0)\mid S=1\bigr].
\]
The loss from diverting the group is $L=\E[Y\mid S=1]-\theta$; its first term is observed directly. When units generate several observations, $P$, $Q$ and $\theta$ are taken at the observation level (Section~\ref{sec:inference}).

\paragraph{Regressions.} Let $m_0$ be a version of $\E[Y\mid X,D=0]$, that is, an $\A$-measurable, $Q$-integrable function such that, for all $A\in\A$,
\[
\E\bigl[Y\,\ind\{X\in A\}\mid D=0\bigr]=\int_A m_0\dd Q .
\]
It exists by the Radon--Nikodym theorem applied to the measure $A\mapsto\E[Y\ind\{X\in A\}\mid D=0]$, which is absolutely continuous with respect to $Q$, and it is unique $Q$-almost surely. With discrete $X$, $m_0(x)$ is the mean outcome of profile $x$ (a rate when $Y$ is binary); the definition through sets is needed because, with continuous $X$, each single value has probability zero. Similarly, $m_0^S$ is a version of $\E[Y(0)\mid X,S=1]$, unique $P$-almost surely.

\paragraph{Assumptions.}
\begin{enumerate}[label=(A\arabic*)]
  \item \emph{Consistency.} $Y=Y(0)$ on $\{D=0\}$. It fails if the basic channel changes when it receives the diverted units, for instance through congestion.
  \item \emph{Conditional mean exchangeability.} $m_0^S=m_0$ $P$-almost surely, where, by (A1), $m_0$ is also a version of $\E[Y(0)\mid X,D=0]$. It concerns only the basic channel: among units with the same profile, those in the target group would obtain there, on average, the same outcome as the units the basic channel serves today. It says nothing about $Y(1)$ and cannot be tested with the data.
  \item \emph{Support.} $P\ll Q$: every set of profiles with probability zero in the donor channel has probability zero in the target group.
  \item \emph{Integrability.} $\E[|Y(0)|\mid S=1]<\infty$ and $\E[|Y|\mid D=0]<\infty$; it holds if $Y$ is bounded.
\end{enumerate}

\begin{remark}[A2 presupposes A3]\label{rem:a2a3}
Since $m_0$ is determined only $Q$-almost surely, the $P$-almost-sure equality in (A2) depends on the chosen version unless $P\ll Q$. If $P$ has mass where $Q$ has none, $m_0$ is arbitrary there and no donor data pin it down. This reappears in Theorem~\ref{thm:lebesgue}: the unsupported part of the target group is not estimable, not merely hard to estimate.
\end{remark}

\begin{remark}[Structural decomposition of the loss]\label{rem:decomposition}
If the outcome is a sum of components, $Y=\sum_kY_k$ (for example, purchases of different products), and the basic channel cannot produce the components in a set $K_{\mathrm{ex}}$, that is, $Y_k(0)=0$ almost surely for $k\in K_{\mathrm{ex}}$, then, by linearity,
\[
L=\sum_{k\in K_{\mathrm{ex}}}\E[Y_k\mid S=1]+\sum_{k\notin K_{\mathrm{ex}}}\bigl(\E[Y_k\mid S=1]-\E[Y_k(0)\mid S=1]\bigr).
\]
Reading $L$ as the causal loss requires consistency also in the premium channel ($Y=Y(1)$ on $\{D=1\}$). The first sum is observed without any identification assumption; only the second needs Sections~\ref{sec:identification} and~\ref{sec:classifier}. An error $\epsilon$ in $\theta_{\mathrm{com}}=\sum_{k\notin K_{\mathrm{ex}}}\E[Y_k(0)\mid S=1]$ produces an error $-\epsilon$ in $\hat L$, so the conclusion about $L$ is insensitive to failures of (A2) or (A3) when $\theta_{\mathrm{com}}/L$ is small. A large structural part is not enough on its own: with $a_{\mathrm{ex}}$ and $a_{\mathrm{com}}$ the two sums of $\E[Y_k\mid S=1]$ and $\rho=a_{\mathrm{ex}}/L$, we have $\theta_{\mathrm{com}}/L=a_{\mathrm{com}}/L-(1-\rho)$, which can be large even when $\rho$ is close to 1. Two caveats: if the unit later obtains the component through another route, $L$ is an upper bound on the loss to the organization; and with a binary ``buys anything'' outcome the decomposition is not additive, so one must work with the component indicators. The full swap, which also moves a group from the basic channel into the premium one, is the same construction with the roles of the channels exchanged.
\end{remark}
